% =============================================================================
\section{Game-Theoretic Analysis}
\label{sec:game-theory}
% =============================================================================

This section provides formal game-theoretic analysis of the PURECORTEX
mechanism across six strategic dimensions: the stake-vs-sell decision, repeated
game dynamics, creator commitment, bonding curve signaling, governance
equilibrium, and delegation rationality.

\subsection{Stake vs.\ Sell Game}
\label{subsec:stake-sell}

Consider a symmetric two-player game in which each player $i \in \{1, 2\}$
holds CORTEX tokens and chooses between Stake (S) and Sell (L). Let:
\begin{itemize}
  \item $e > 0$: per-period staking emission reward when both players stake.
  \item $g > 0$: capital appreciation from (a) reduced circulating supply via
    locking and (b) continuous buyback-burn via the Assistance Fund.
  \item $c > 0$: opportunity cost of locking capital (illiquidity penalty).
  \item $\ell > 0$: immediate liquidation value.
  \item $d > 0$: price depreciation suffered by the remaining staker when the
    other sells.
\end{itemize}

The payoff matrix is:

\begin{table}[H]
\centering
\caption{Stake vs.\ Sell payoff matrix.}
\label{tab:stake-sell}
\begin{tabular}{@{}lcc@{}}
\toprule
& \textbf{Player 2: Stake} & \textbf{Player 2: Sell} \\
\midrule
\textbf{Player 1: Stake} & $(e + g - c,\; e + g - c)$ & $(e - d - c,\; \ell)$ \\
\textbf{Player 1: Sell}  & $(\ell,\; e - d - c)$ & $(\ell - d,\; \ell - d)$ \\
\bottomrule
\end{tabular}
\end{table}

\begin{proposition}[Nash Equilibrium in Stake-Sell Game]
\label{prop:stake-nash}
The strategy profile $(\text{Stake}, \text{Stake})$ is a Nash equilibrium if
and only if:
\begin{equation}
  e + g - c \geq \ell
  \label{eq:stake-condition}
\end{equation}
That is, the combined staking emission rewards and capital appreciation
(from both locking-driven supply reduction and Assistance Fund buyback-burn)
exceed the liquidation value plus the opportunity cost.
\end{proposition}

\begin{proof}
For (Stake, Stake) to be a Nash equilibrium, neither player can profitably
deviate. Player 1's payoff from Stake (given Player 2 stakes) is $e + g - c$.
Player 1's payoff from deviating to Sell is $\ell$. The no-deviation condition
is $e + g - c \geq \ell$, which is condition~\eqref{eq:stake-condition}.
By symmetry, the same condition applies to Player 2. \qed
\end{proof}

\begin{remark}
The \veCORTEX{} mechanism is designed to ensure condition~\eqref{eq:stake-condition}
holds through three reinforcing channels:
\begin{enumerate}
  \item \textbf{Emission rewards ($e$):} The boost multiplier increases
    emission yield for longer lock durations, funded by the 24\% staking pool.
  \item \textbf{Buyback-burn appreciation ($g$):} The Assistance Fund
    continuously buys and burns CORTEX using 90\% of all protocol revenue,
    creating sustained upward price pressure that benefits all holders but
    disproportionately rewards stakers (who cannot sell into the appreciation).
  \item \textbf{Governance power:} Stakers control protocol parameters via
    \veCORTEX{}-weighted voting, including the ability to adjust fee rates and
    emission schedules.
\end{enumerate}
The lock duration commitment ensures that the ``sell'' option is not available
during the lock period, effectively removing the deviation possibility for
committed stakers.
\end{remark}

\subsection{Repeated Game and the Folk Theorem}
\label{subsec:repeated-game}

The one-shot stake-sell game may admit (Sell, Sell) as an equilibrium when
$e + g - c < \ell$ (e.g., during bear markets). However, the \veCORTEX{} lock
mechanism converts this one-shot game into a repeated game.

\begin{theorem}[Cooperation via Lock Commitment]
\label{thm:folk}
Let the stake-sell game be repeated $T$ times, where $T = T_{\text{lock}} /
\Delta t$ is the number of periods in the lock duration. By the folk
theorem~\citep{fudenberg1991game}, if the discount factor $\delta$ satisfies:
\begin{equation}
  \delta \geq \frac{\ell - (e + g - c)}{(\ell - d) - (e - d - c)} = \frac{\ell - e - g + c}{\ell - e + c}
  \label{eq:folk-threshold}
\end{equation}
then the cooperative outcome $(\text{Stake}, \text{Stake})$ can be sustained
as a subgame-perfect equilibrium using grim-trigger strategies.
\end{theorem}

The proof follows from the standard folk theorem argument and is provided in
Appendix~\ref{app:proofs}. The key insight is that the \veCORTEX{} lock
literally removes the ability to deviate for the lock duration, making
cooperation not merely incentive-compatible but mechanically enforced. This is
stronger than the folk theorem's requirement of voluntary cooperation sustained
by threat of punishment.

\subsection{Creator Commitment}
\label{subsec:creator-commitment}

We model the creator's optimal strategy using a Bellman equation. Let $W(d)$
denote the creator's expected wealth at day $d$ of the vesting schedule:

\begin{equation}
  W(d) = V(d) \cdot P(d) + \sum_{t=d}^{\infty} \beta^{t-d} \cdot f_{\text{creator}} \cdot TV(t)
  \label{eq:bellman-creator}
\end{equation}

where:
\begin{itemize}
  \item $V(d)$ is the vested token quantity (Equation~\eqref{eq:creator-vesting}).
  \item $P(d)$ is the token price at day $d$.
  \item $\beta \in (0, 1)$ is the creator's discount factor.
  \item $f_{\text{creator}}$ is the creator's revenue share from agent-specific
    activity (agent token trading, tool invocations).
  \item $TV(t)$ is the trading volume at time $t$.
\end{itemize}

A creator who sells all vested tokens at day $d$ realizes $V(d) \cdot P(d)$
but forfeits all future fee revenue. A creator who holds (and continues
developing the agent) maintains the fee stream but bears token price risk.

\begin{proposition}[Vesting Maximizes Expected Wealth]
\label{prop:creator-wealth}
Under the assumption that creator effort $e(d)$ positively affects both
$P(d)$ and $TV(t)$ for $t > d$, the vesting schedule maximizes the creator's
expected lifetime wealth compared to immediate full allocation, because:
\begin{equation}
  \frac{\partial W}{\partial e} = V(d) \cdot \frac{\partial P}{\partial e} + \sum_{t=d}^{\infty} \beta^{t-d} \cdot f_{\text{creator}} \cdot \frac{\partial TV}{\partial e} > 0
  \label{eq:effort-incentive}
\end{equation}
The vesting constraint forces the creator to internalize the long-term
consequences of effort withdrawal, since immediate liquidation of the full
allocation is impossible.
\end{proposition}

The proof is provided in Appendix~\ref{app:proofs}. The intuition is
straightforward: by restricting the creator's ability to sell in the short term,
the vesting schedule converts the creator's optimization problem from a
short-horizon liquidation problem to a long-horizon wealth-maximization problem,
aligning incentives with ecosystem health.

\subsection{Bonding Curve as Quality Signal}
\label{subsec:signaling}

We adapt the Spence signaling model~\citep{spence1973job} to the agent
tokenization context.

\begin{definition}[Agent Quality Model]
\label{def:quality-model}
An agent creator has private information about agent quality
$\theta \in \{\theta_L, \theta_H\}$, where $\theta_H > \theta_L$. The
creator chooses to launch on PURECORTEX (quadratic bonding curve with
vesting) or on a platform with no vesting (flat bonding curve). Buyers
observe the platform choice but not $\theta$ directly.
\end{definition}

\begin{theorem}[Separating Equilibrium via Quadratic Curve]
\label{thm:separating}
Under the quadratic bonding curve with 180-day creator vesting, a separating
equilibrium exists in which:
\begin{enumerate}
  \item[(i)] High-quality creators ($\theta_H$) choose PURECORTEX.
  \item[(ii)] Low-quality creators ($\theta_L$) choose the no-vesting platform.
  \item[(iii)] Buyers correctly infer quality from platform choice.
\end{enumerate}
The separating condition requires:
\begin{equation}
  c_H(\text{vest}) < \theta_H \cdot \Delta P < c_L(\text{vest})
  \label{eq:separating}
\end{equation}
where $c_\theta(\text{vest})$ is the cost of the vesting constraint for a
creator of type $\theta$, and $\Delta P$ is the price premium that buyers
assign to PURECORTEX-launched agents.
\end{theorem}

The proof is provided in Appendix~\ref{app:proofs}. The key mechanism is that
the 180-day vesting imposes a higher effective cost on low-quality creators
(who cannot sustain agent development and thus face declining token prices
during the vesting period) than on high-quality creators (whose continued
development maintains or increases token prices). This cost differential
enables quality separation.

\subsection{Senator Mechanism Design}
\label{subsec:senator-mechanism}

\begin{theorem}[Pareto Improvement via Governance]
\label{thm:pareto}
Let $\mathcal{O} = \{o_1, \ldots, o_K\}$ be the set of possible protocol
parameter configurations, and let $u_i(o)$ be agent $i$'s utility under
configuration $o$. If the Senator maximizes the welfare function $W$ (Equation~\eqref{eq:welfare}) and the governance mechanism selects the
configuration $o^*$ that maximizes $\sum_i u_i(o)$ subject to constitutional
constraints, then $o^*$ is a Pareto improvement over the status quo $o_0$
whenever:
\begin{equation}
  W(o^*) > W(o_0) + \epsilon
  \label{eq:pareto-condition}
\end{equation}
for sufficiently small $\epsilon > 0$ (the proposal significance threshold).
\end{theorem}

\begin{proof}
The welfare function $W$ is a weighted sum of metrics (TVL, agent count,
retention, negative risk) that are positively correlated with individual agent
utilities. An increase in TVL benefits stakers through higher revenue share; an
increase in agent count benefits all agents through composability network
effects; an increase in retention reflects revealed preference for the current
configuration; and a decrease in risk benefits all participants.

Formally, for each component $W_k$ of the welfare function:
\begin{equation}
  \frac{\partial u_i}{\partial W_k} \geq 0 \quad \forall i, k
  \label{eq:welfare-monotone}
\end{equation}

If $W(o^*) > W(o_0)$, then at least one component $W_k$ has increased, and no
component has decreased (since the Senator only proposes changes that increase
$W$). By the monotonicity of individual utilities in welfare components, at
least one agent is strictly better off and no agent is worse off---the
definition of a Pareto improvement. \qed
\end{proof}

\subsection{Delegation Game}
\label{subsec:delegation-game}

\begin{proposition}[Rationality of Delegation]
\label{prop:delegation}
Delegation dominates direct voting for agent $i$ when the information
acquisition cost exceeds the voting cost:
\begin{equation}
  c_{\text{info}}^i > c_{\text{vote}}^i + (VP_i^{\text{eff}} - \veCORTEX_i) \cdot \frac{\partial u_i}{\partial \text{outcome}}
  \label{eq:delegation-condition}
\end{equation}
where $c_{\text{info}}^i$ is the cost of acquiring sufficient information to
vote optimally, $c_{\text{vote}}^i$ is the transaction cost of voting, and the
third term captures the influence premium from receiving delegations (which is
zero for a pure delegator).
\end{proposition}

\begin{proof}
An agent choosing between direct voting and delegation compares expected
utilities. Under direct voting, the agent pays $c_{\text{info}}^i +
c_{\text{vote}}^i$ and casts a vote that may or may not be well-informed.
Under delegation to a trusted delegate $j$ with superior information, the
agent pays zero cost and the delegate votes on their behalf with higher
expected quality.

Let $p_i$ be the probability that agent $i$ votes correctly when investing
$c_{\text{info}}^i$ in information acquisition, and let $p_j > p_i$ be the
delegate's probability of voting correctly (reflecting the delegate's
comparative advantage in governance analysis).

The expected utility from direct voting is:
\begin{equation}
  EU_{\text{direct}}^i = p_i \cdot u_i(\text{correct}) + (1 - p_i) \cdot u_i(\text{incorrect}) - c_{\text{info}}^i - c_{\text{vote}}^i
  \label{eq:eu-direct}
\end{equation}

The expected utility from delegation is:
\begin{equation}
  EU_{\text{delegate}}^i = p_j \cdot u_i(\text{correct}) + (1 - p_j) \cdot u_i(\text{incorrect})
  \label{eq:eu-delegate}
\end{equation}

Delegation dominates when $EU_{\text{delegate}}^i > EU_{\text{direct}}^i$:
\begin{equation}
  (p_j - p_i) \cdot [u_i(\text{correct}) - u_i(\text{incorrect})] > -c_{\text{info}}^i - c_{\text{vote}}^i
  \label{eq:delegation-dominance}
\end{equation}

Since $p_j > p_i$ and the utility difference is positive, the left side is
positive, and the condition is satisfied whenever the information and voting
costs are non-negative---which they always are. More precisely, delegation
dominates whenever the delegate's information advantage outweighs any
agency costs. \qed
\end{proof}

\subsection{Constitutional Rigidity as Schelling Point}
\label{subsec:schelling}

The supermajority requirement for constitutional amendments (67\%) serves a
game-theoretic function beyond simple conservatism. In a voting game where
participants must coordinate on which rules to preserve, the supermajority
threshold creates a Schelling point~\citep{fudenberg1991game}:

\begin{itemize}
  \item A simple majority (50\% + 1) allows a minimal winning coalition to
    impose changes that benefit the majority at the expense of the minority
    (tyranny of the majority).
  \item A supermajority (67\%) requires broad consensus, ensuring that
    constitutional changes reflect near-unanimous agreement on the direction
    of protocol evolution.
  \item The 67\% threshold is strategically chosen: it is high enough to
    prevent narrow coalitions from capturing the governance process, but low
    enough to permit constitutional evolution when circumstances genuinely
    warrant it.
\end{itemize}

This creates a credible commitment device: participants can trust that
foundational rules will not be changed by a narrow interest group, encouraging
long-term investment in the ecosystem.
